\section{Auditoría de fuentes y límites del problema: esta es una clase de alineación de 2023}

Esta sesión se grabó el 17 de enero de 2023, poco después de que ChatGPT se hiciera público. Trata la ruta de RLHF que ya estaba consolidada en ese momento, la scalable oversight que se estaba explorando y las preguntas abiertas que plantearon los estudiantes sobre preferencias, engaño, herramientas, inner alignment e interpretabilidad. La versión anterior trasladaba a la clase hechos posteriores: el relato organizacional de ``Superalignment'', procesos genéricos de puesta en producción y listas de gobernanza de producto. Por eso, esta vez se reescribió por completo con base en el video oficial, 1,264 subtítulos hechos por personas, 17 imágenes didácticas de la clase y seis artículos de primera mano ya publicados en ese momento, en lugar de ampliar la estructura anterior a base de parches.

\lecturefigure{slide-01-title.jpg}{Título de la sesión de Jan Leike: Aligning Language Models with Humans.}{Video oficial de Stanford Online 00:00:23--00:00:49.}

En el título, ``with humans'' es más concreto que ``make models safe'': la sesión pregunta cómo pueden las personas convertir en señal de entrenamiento intenciones que no se pueden escribir por completo como reglas, y cómo mantener útil esa señal mientras la capacidad de los modelos sigue aumentando. El lector debe distinguir tres tipos de evidencia: las slides presentan la línea pública de la exposición, las explicaciones orales del ponente añaden motivaciones y limitaciones, y los artículos de primera mano aportan fórmulas y detalles experimentales. Cuando las tres no coinciden, la nota marca el límite de forma explícita en lugar de hacer pasar conocimiento posterior por palabras textuales de la clase.

\begin{importantbox}{El contrato de evidencia de esta sesión}
Las cifras, el estado de los sistemas y los juicios de investigación que aparecen en la clase tienen como límite el 2023-01-17. El artículo público de InstructGPT puede usarse para aclarar la relación entre SFT, reward model y PPO; los nombres de proyectos, las arquitecturas de despliegue y los cambios de política posteriores a 2023 no pueden presentarse como respuestas que el ponente ya había dado entonces. Durante la última media hora se muestra casi todo el tiempo la última slide, por lo que la teacher voice proviene principalmente de los subtítulos con marca de tiempo, y no de insertar repetidamente la misma imagen del ponente.
\end{importantbox}

\subsection{Team AI: por qué el aumento de capacidades convierte la alineación en un problema estratégico}

El ponente no empieza con funciones de pérdida, sino que plantea una competencia entre ``Team AI y Team Human''. El objetivo de esta analogía no es afirmar que los sistemas actuales ya superan a las personas en todo, sino introducir la dinámica en el problema: distintos AI player van entrando poco a poco, y los player posteriores pueden ser más rápidos, más baratos y más capaces en más tareas. Por ahora, las personas conservan una ventaja importante: decidir qué sistemas entran al mundo real y de qué manera, y entrenar a algunos de ellos para que se pongan del lado humano.

\lecturefigure{slide-02-team-ai-strong-players.jpg}{Team AI se está incorporando a la competencia, y en el futuro podrían aparecer player sumamente fuertes.}{Video oficial de Stanford Online 00:00:49--00:03:18.}

\begin{knowledgebox}{Lectura de la figura: un player fuerte no significa «hoy ya lo puede todo»}
Esta página ve la capacidad de los modelos como una secuencia de player que se incorporan con el tiempo. El ponente subraya a la vez que los sistemas de 2023 todavía tienen limitaciones evidentes; ChatGPT es solo un punto de referencia que muestra cómo la amplitud de conocimientos, la cobertura de idiomas y un costo marginal bajo podrían cambiar la división del trabajo. El ``incredibly strong'' de la figura es una hipótesis de modelado de riesgos orientada al futuro, no una conclusión medida sobre cualquier modelo contemporáneo.
\end{knowledgebox}

Después, la analogía se divide en dos objetivos de acción. Primero, entrenar o elegir AI player dispuestos a ayudar a las personas y a actuar conforme a la intención humana; segundo, diseñar instituciones y reglas para que un entorno competitivo con IA poderosa no lleve a la humanidad a la derrota. El ponente llama alignment al primero; el segundo se acerca más a governance. Ambos se influyen mutuamente, pero no son el mismo problema técnico.

\lecturefigure{slide-03-two-objectives.jpg}{Los dos objetivos de Team Human: reclutar AI player y establecer reglas con las que no se pierda la competencia.}{Video oficial de Stanford Online 00:03:18--00:04:16.}

\begin{warningbox}{No mezclar alignment y governance en un término que lo abarque todo}
Alignment estudia señales de entrenamiento, objetivos de comportamiento, evaluación y generalización; governance se ocupa de los permisos de despliegue, las reglas de competencia, la responsabilidad y la coordinación institucional. Esta sesión solo desarrolla lo primero. Introducir en lo que sigue procesos organizacionales, propuestas regulatorias o la competencia de mercado no solo se aparta de la clase, sino que también oculta el verdadero cuello de botella técnico de la sesión: cuando las personas no entienden las tareas que realiza el modelo, cómo seguir dando retroalimentación confiable.
\end{warningbox}

\subsection{Resumen de la sección}

Esta sección establece dos límites. En el tiempo, se trata de una clase de principios de 2023; en el problema, se centra en cómo lograr que los modelos sigan la human intent, sin intentar resolver toda la AI governance. La analogía de Team AI aporta la motivación estratégica; lo siguiente es descomponer ``ponerse del lado humano'' para que deje de ser un lema y se convierta en objetivos entrenables.
